% Whole-family criterion; full root and representations independent written audits passed.
\begin{theorem}[Proved: the complete criterion for all unramified dyadic normal nodes]
\label{C:dyadic_normal_node_criterion}
Let $A$ be any characteristic-zero DVR with maximal ideal $(2)$,
let $\kappa=A/2A$, let $u\in A^\times$, and let $n\ge1$.
Take either of the two local rings
\[
 R_n=\left(A[x,y]/(xy-2^nu)\right)_{(2,x,y)}
 \quad\hbox{or}\quad
 R_n^{\mathrm{for}}=A[[x,y]]/(xy-2^nu).                    \tag{C.1.N7}
\]
Both are flat, normal, two-dimensional noetherian local domains
over $A$. For either choice $R$, the following are equivalent:
\begin{enumerate}
\item $\vartheta_R$ is an isomorphism in every degree;
\item its degree-two map is surjective;
\item $n=1$;
\item $R$ is regular.
\end{enumerate}
For every $n\ge2$ there is an explicit nonzero class in
\[
       \ker\bigl(R_R(C_2)\longrightarrow R_{R[1/2]}(C_2)\bigr).
                                                               \tag{C.1.N8}
\]
The detecting representation has rank two for $n\ge3$ and rank
six for $n=2$, independently of the residue field.
For the polynomial localization $R_n$, every nonzero
integer localization has $K_0=\mathbb Z$: its generic ring
is a localization of $\operatorname{Frac}(A)[x,x^{-1}]$.
No generic $K_0$ assertion for the formal ring is needed here.
\end{theorem}
\begin{proof}
\emph{The rings and the positive case.}
The polynomial ring has the domain, dimension, flatness and
normality proof of Theorem~\ref{C:normal_trivialK0_kernel},
with $4$ replaced by $2^nu$. The coefficient polynomial is
primitive and irreducible over $A[x]$. Its generic ring is
the stated localization of a Laurent polynomial PID.
The only vertical height-one primes are $(2,x)$ and $(2,y)$;
at each the other variable is a unit and the maximal ideal
is generated by two. The hypersurface satisfies $(S_2)$,
so the verified $(R_1)$ condition gives normality.
Locality and the generic PID give the $K_0$ assertions exactly
as before.

For completeness, the formal ring has the same properties
without assuming an identification of its generic Picard group.
The ring $S=A[[x,y]]$ is a regular local ring of dimension three:
its maximal ideal has generators $2,x,y$, and the evident chain
of primes has length three. It is a UFD
\cite[Lemma 15.123.2, Tag 0AG0]{C-NormalNodesStacks}.
We check that $F=xy-2^nu$ is irreducible in $S$.
If $F=gh$ with both factors nonunits, reduction modulo two
in the UFD $\kappa[[x,y]]$ shows, after interchanging them,
that $\bar g=xv$ and $\bar h=yw$ with $v,w$ units.
Setting $x=0$ gives
\[
                  g(0,y)h(0,y)=-2^nu.
\]
The equality holds over the UFD $A[[y]]$. But $h(0,y)$ is neither a unit nor
divisible by two, because its reduction is $yw(0,y)$.
This contradicts the displayed factorization of a unit times
a power of the prime two. Thus $F$ is prime, and the quotient
is a domain of dimension two, flat over $A$.

It is again a Cohen--Macaulay hypersurface. The vertical
height-one localizations are DVRs as above. For any other
prime $\mathfrak p$ of the quotient, let $\mathfrak q$ be its
preimage in $S$. If two is invertible there, then $x$ and $y$
are units modulo $F$, and hence $x\notin\mathfrak q$.
The formal derivation $\partial/\partial y$ extends to
$S_{\mathfrak q}$ and sends $F$ to $x$.
Were $F\in\mathfrak q^2S_{\mathfrak q}$, its derivative
would belong to $\mathfrak qS_{\mathfrak q}$, a contradiction.
A quotient of a regular local ring by an element outside
the square of its maximal ideal is regular. Thus $(R_1)$
holds also here, and Serre's criterion gives normality.

When $n=1$, the equation has a nonzero linear term in two
in the ambient cotangent space, so either quotient is regular.
When $n\ge2$, the equation is in the square of the ambient
maximal ideal, so the quotient has embedding dimension three
and dimension two and is singular. This proves (3)$\Leftrightarrow$(4).
For $n=1$, Theorem~\ref{C:regular_noetherian} gives (1);
(1)$\Rightarrow$(2) is immediate. We now disprove (2) for
every $n\ge2$, with no hypothesis on $\kappa$.

\emph{A common residue detector.}
If $A'$ is a DVR extension of $A$ with maximal ideal $(2)$,
put
\[
                 D=A'[s]_{(2)},\qquad
                 D/2D=\kappa'(s),\quad \kappa'=A'/2A'.  \tag{C.1.N9}
\]
This is a DVR with uniformizer two: $A'[s]$ is a domain,
$(2)$ is prime, and its localization has that principal maximal
ideal. The field parameter $s$ is transcendental over $\kappa'$.
For every $n\ge2$ there is a ring map
\[
                    R\longrightarrow D/4D,
                         \qquad x\mapsto2,\quad y\mapsto2s.
                                                               \tag{C.1.N10}
\]
It respects the equation since both $4s$ and $2^nu$ vanish.
For the polynomial localization every inverted denominator has
nonzero residue in $\kappa$, hence is a unit in the target.
For the formal ring evaluation is well-defined because the images
of $x,y$ generate a square-zero ideal; only their constant and
linear terms contribute. This also proves that the map respects
all formal power-series relations.

For any field $F$ of characteristic two and any nonzero $c\in F$,
a polynomial
\[
                       z^2+z+(cs+d),\qquad d\in F,
\]
is irreducible over $F(s)$. Indeed $cs+d$ has a pole of order
one at infinity. If $w\in F(s)$ has a pole there, $w^2+w$
has a pole of even order; if $w$ has no pole there, neither
does $w^2+w$. Thus the polynomial has no root in $F(s)$.
It is invariant under $z\mapsto z+1$.
Lemma~\ref{C:dyadic_residue_Bockstein} therefore applies to
all such irreducible primaries over $D/4D$.

\emph{Every exponent $n\ge3$.}
Here take $A'=A$ and put
\[
 a=1+2^{n-1},\qquad v=u^{-1}(1+2^{n-2}),\qquad
 T=\begin{pmatrix}a&x\\-vy&-a\end{pmatrix}.             \tag{C.1.N11}
\]
The identity $a^2-1=2^n(1+2^{n-2})$ gives $T^2=I_2$.
After the map (C.1.N10), the reduced differential is zero and
its Bockstein is
\[
             N=\begin{pmatrix}0&1\\\bar u^{-1}s&1\end{pmatrix}.
\]
Its characteristic polynomial is $z^2+z+\bar u^{-1}s$,
which is irreducible by the preceding pole argument.
The corresponding primary has length one, so its parity is one.
The trivial and sign lines have parity zero.
It follows that $[T]-1-\mathrm{sgn}$ is nonzero and lies outside
the degree-two polynomial subgroup.

\emph{The remaining exponent $n=2$.}
Let $C$ be an integral companion matrix of
$q(t)=t^3+t+1$ and set $a=I_3+2C$. Define
\[
 T=\begin{pmatrix}
       a&xI_3\\
       -u^{-1}y(C^2+C)&-a
    \end{pmatrix}.                                      \tag{C.1.N12}
\]
All its blocks commute, and
$a^2-I_3=4(C^2+C)$. The node equation again gives $T^2=I_6$.

We construct the required $A'$ without a perfectness or
henselianity assumption on $A$ or $\kappa$.
Over $\mathbb F_2$, $q$ is irreducible, and its three roots
are $\beta,\beta^2,\beta^4$ in $\mathbb F_8$.
Over any field $\kappa$ of characteristic two, either it has
no root and is irreducible of degree three, or one root supplies
all three of these roots in $\kappa$.
In the split case put $A'=A$.
In the irreducible case put
$A'=A[b]/(b^3+b+1)$.
This is a finite free local domain: reduction modulo two is
a field, and Gauss's lemma over the DVR makes the monic
polynomial irreducible over $\operatorname{Frac}(A)$.
Its maximal ideal is generated by two, so it is a DVR.
Its defining derivative is nonzero at every residue root,
so it is also unramified (indeed finite \'etale) over $A$.
Thus in both cases the residue field $\kappa'$ splits $q$.

Write its roots as $\beta_1,\beta_2,\beta_3$ and put
$d_i=\beta_i^2+\beta_i$. These three elements are nonzero
and pairwise distinct. Nonzero follows because neither zero
nor one is a root of $q$. If $d_i=d_j$ for $i\ne j$, then
$\beta_i+\beta_j=1$; the zero coefficient of $t^2$ in $q$
would make the third root equal to one, a contradiction.

Over $\kappa'$ the reduction of $C$ is diagonalizable with
these three distinct eigenvalues. After (C.1.N10) the
Bockstein therefore decomposes over $\kappa'(s)$ into the
three matrices
\[
 N_i=\begin{pmatrix}
        \beta_i&1\\
        \bar u^{-1}s d_i&1+\beta_i
     \end{pmatrix},\qquad
 f_i(z)=\det(zI-N_i)=z^2+z+d_i(1+\bar u^{-1}s).
\]
Each $f_i$ is irreducible by the pole argument, and the three
polynomials are distinct because the $d_i$ are distinct.
For $f=f_1$, the primary space has dimension two and length
one. The parity invariant is again one, whereas it vanishes
on both constant sign lines.
Consequently $[T]-3(1+\mathrm{sgn})$ is nonzero and is outside
the degree-two polynomial subgroup.

\emph{Generic vanishing and the equivalence.}
For either displayed matrix write its size as $2d$ with
diagonal block $a$ ($d=1$ or $3$). After inverting two,
$x$ is a unit, and the block matrix
\[
       P=\begin{pmatrix}xI_d&xI_d\\I_d-a&-I_d-a\end{pmatrix}
\]
has determinant $(-2x)^d$ and obeys
$TP=P\operatorname{diag}(I_d,-I_d)$.
Thus both generic eigenspaces are explicitly free, and the
nonzero class just detected belongs to (C.1.N8).
The degree-two polynomial image is the span of the trivial
and sign lines, by the standard admissible regular-module
sequence. The detector annihilates that image and takes
value one on the constructed representation. Hence (2)
fails whenever $n\ge2$, completing all equivalences.
\end{proof}

The theorem in particular settles nonvanishing of the degree-two kernel for
$\mathbb Z_2[[x,y]]/(xy-4)$ itself. Its proof uses a
transcendental residue-field test after a finite unramified
extension; it does not assume that the original residue field
has an Artin--Schreier extension. The criterion is exactly
regularity within the family (C.1.N7), not a regularity
criterion for arbitrary mixed-characteristic rings.

A particularly small example over the original dyadic base is
\[
 R=\left(\mathbb Z_{(2)}[x,y]/(xy+8)\right)_{(2,x,y)},
 \qquad T=\begin{pmatrix}3&x\\y&-3\end{pmatrix}.
\]
Here $T^2=I_2$, and $x,y\mapsto2$ gives the Bockstein
$\left(\begin{smallmatrix}1&1\\1&0\end{smallmatrix}\right)$
over $\mathbb F_2$, detected by $t^2+t+1$.
Thus this single normal domain has trivial $K_0$ at every
integer localization and a nonzero generic kernel without
any coefficient-field extension.

The bounded exact certificate
\texttt{code/C\_finish\_normal\_nodes.py} checks the symbolic
involution and generic-eigenbasis identities for $2\le n\le10$,
and the three distinct primary blocks over $\mathbb F_8[s]$
for every nonzero residue unit in $\mathbb F_8$.
It also checks the last $xy+8$ example. Its output and runtime
are in \texttt{logs/C\_finish\_normal\_nodes.json}.
The general exponent, arbitrary residue field, pole argument
and commutative-algebra assertions are established by the
written proof, beyond that bounded certificate.
